\section{Introduction}
Egyptian fruit bats (\emph{Rousettus aegyptiacus}) are social, vocal mammals with a large public annotated call corpus \cite{prat2017} whose calls carry emitter, addressee, context and behaviour information \cite{prat2016}. Whether their acoustic space holds discrete call types or varies continuously matters for describing repertoires: hard clustering handles gradation poorly and fuzzy clustering has been recommended \cite{wadewitz2015,fischer2017}, while latent-space projection with density clustering \cite{sainburg2020,thomas2022} and autoencoder or pretrained embeddings \cite{goffinet2021,best2023,aves2023,stowell2022} are now widely used.

We audit two earlier analyses of this corpus (Section~2). (Q1) Can the original analysis be reproduced, and do its labels belong to its audio? (Q2) Do its clusters follow silence, duration, noise, channel, session or date? (Q3) Is there discrete structure in voiced frames and syllables when one cluster is allowed, individuals are resampled and power is known? (Q4) Does any call group survive validation blocked by individual, channel and session, and replicate in unused recordings? (Q5) How much individual and context information transfers across recording conditions? Out-of-sample criteria were fixed beforehand, and part of the original analysis rested on labels that did not belong to its audio.

\section{Background}
Such analyses are fragile for four reasons. First, fuzzy c-means \cite{bezdek1984} and its validity indices \cite{xiebeni1991,rousseeuw1987} cannot return a single cluster, clusterability tests \cite{adolfsson2019,liu2008} are rarely reported, and low-dimensional maps can mislead \cite{chari2023}. Second, clusters may follow recording properties (silence, background noise, microphone) rather than call structure. Third, calls from one file, individual or session are not independent, so ordinary tests and random splits overstate evidence \cite{hurlbert1984} and generalisation \cite{roberts2017}. Fourth, in colony recordings individuals are confounded with place and time.

Two earlier analyses applied FCM to summary features of this corpus: a June analysis described ``a fluid, continuous manifold'' with individual ``acoustic modes'', and a July manuscript reported two ``reproducible acoustic regimes'' linked to context and emitter.
